\documentclass[11pt,a4paper]{article}
\usepackage[utf8]{inputenc}
\usepackage[T1]{fontenc}
\usepackage{lmodern}
\usepackage[margin=1in]{geometry}
\usepackage{amsmath,amssymb,amsthm}
\usepackage{booktabs}
\usepackage{array}
\usepackage{microtype}
\usepackage[hidelinks]{hyperref}

\newtheorem{theorem}{Theorem}
\newtheorem{proposition}[theorem]{Proposition}
\newtheorem{lemma}[theorem]{Lemma}
\theoremstyle{definition}
\newtheorem{observation}[theorem]{Observation}
\theoremstyle{remark}
\newtheorem{remark}[theorem]{Remark}

\newcommand{\tier}[1]{{\small\textsf{[#1]}}}
\newcommand{\lmin}{\lambda_{\min}}
\newcommand{\src}[1]{{\footnotesize\texttt{#1}}}
\newcommand{\Qrest}{Q_{\mathrm{rest}}}

\title{\bfseries A Certified Negative Finite-Window Weil Witness for the
Davenport--Heilbronn Function, with the Remaining-Zero Form Positive}

\author{Elias De Jes\'us\\[2pt]
{\small Independent Researcher}\\[1pt]
{\small ORCID \href{https://orcid.org/0009-0007-0190-9143}{0009-0007-0190-9143}}\\[1pt]
{\small \texttt{dejesuselias10@gmail.com}}}

\date{October 2026 --- Technical Note\\[2pt]
{\small\texttt{doi:10.5281/zenodo.XXXXXXXX}}}

\begin{document}
\maketitle

\begin{abstract}
\noindent
The Davenport--Heilbronn function $f$ satisfies a $\zeta$-shaped functional equation, is
entire after completion, has all its zeros in an explicit strip, and has no Euler
product; it is known to have zeros off the critical line. It is therefore the natural
negative control for any finite-window Weil positivity detector. We exhibit one
explicitly specified real test function $g_*$ supported in $[-2,2]$, odd, in a
$37$-dimensional span, whose Davenport--Heilbronn Weil value is enclosed in ball
arithmetic as
$Q_f(g_*)/\|g_*\|_2^2 = -0.679071346441480590986377586093\ldots$ with radius
$1.39\times10^{-116}$, and a second such function at $L=1.8$ with value
$-1.68123895523670103\times10^{-16}$. Both are strictly negative. Decomposing against a
census of the zeros to height $200$, the single off-line quartet at
$\beta = 0.808517182456637386$, $\gamma = 85.6993484853775922$ contributes
$-0.985198326439305$, and the entire remaining spectrum contributes $+0.306126979997825$;
the minimum eigenvalue of the remaining-zero form is positive at every scanned window and
sector, collapsing to $8.47\times10^{-36}$ at $L=2.6$ while the full form is at $-2.14$.
On this object, therefore, the hypothesis of the remaining-zero theorem of the companion
paper holds, and the negativity is carried by the off-line quartets alone. What this
rules out is any claim that finite-window Weil positivity follows from a functional
equation, entireness and a zero-containing strip; it says nothing about $\zeta$, and it
cannot separate ``has off-line zeros'' from ``has no Euler product''.
\end{abstract}

\section*{0.\quad Non-claims}
\addcontentsline{toc}{section}{Non-claims}

Stated first, not in an appendix.

\begin{enumerate}\itemsep3pt
\item \textbf{Nothing about RH or GRH.} Every object here is the Davenport--Heilbronn
function. It is not $\zeta$, it is not an $L$-function of anything, and no statement
below transfers to either.

\item \textbf{$f$ has no Euler product, and nothing here transfers to $\zeta$.} A
detector that sees this function's off-line zeros is a detector that is not blind. It is
not thereby a detector that would see $\zeta$'s, if $\zeta$ had any.

\item \textbf{The confound, stated here and not buried.} $f$ differs from $\zeta$ in
\emph{two} ways at once: it has zeros off the critical line, and it has no Euler product.
Every measurement below is consistent with either being the cause of the negativity, and
nothing here separates them. Separating them would need an object with one property and
not the other, and no such object is known to the author. In particular, the observation
of \S5 that the remaining-zero form is positive is evidence \emph{internal to this
object} that the off-line zeros carry the negativity; it is not a proof, because the
remaining-zero form is itself built from a function with no Euler product.

\item \textbf{Tiers.} Every zero location is \tier{T2}. The two witnesses of \S4 are
\tier{T2-C}. The onset scan of \S6 is \tier{T2-N}. \textbf{Non-detections certify
nothing}: a finite basis that finds no negative direction is not evidence of positivity,
and the absence of a witness at $L=1.6$ is reported as an absence, never as positivity.

\item \textbf{The $L=1.8$ onset is not certified.} A negative witness \emph{at} $L=1.8$
is certified (\S4); that the form is positive at $L=1.6$ is not, so the crossing is
bracketed by one certified negative and one non-detection, and the location of the
crossing remains \tier{T2-N}.

\item \textbf{The two literature witnesses are unverified.} Only landing pages were read;
see Appendix~C. No quantity here was fitted to either.

\item \textbf{A finite positive matrix proves positivity only on its own test space.} The
positivity of the remaining-zero form in \S5 is a statement about a $37$-dimensional
restriction at each window, never about the whole finite-window function space.

\item \textbf{Numerical root residuals do not establish census completeness}, and
agreement with a published decimal is a cross-check, not an interval certificate.

\item \textbf{Not treated here.} The degenerate case $\gamma = 0$ of the orbit
contribution, where the parity split behaves differently, is the subject of a companion
note and is not used anywhere below.
\end{enumerate}

\section*{Epistemic tiers}
\addcontentsline{toc}{section}{Epistemic tiers}

\begin{center}
\begin{tabular}{@{}lll@{}}
\toprule
Tier & Meaning & Discharged by \\
\midrule
T1 & exact or proved & a proof, or exact arithmetic over $\mathbb{Z}$ or $\mathbb{Q}$ \\
T2 & numerical & stated working precision $+$ a convergence or noise-floor check \\
T2-C & certified arithmetic & an Arb ball enclosure with a reproducible certificate \\
T3 & interpretation & argument from T1/T2 evidence, labelled as such \\
\bottomrule
\end{tabular}
\end{center}

\noindent T2-C is \emph{not} a formal-proof assistant verification and is nowhere
described as one. Every number printed in this note is re-derived from the data files by
\src{note/verify\_note\_numbers.py}, which also fails if the note prints more digits than
its source carries.

\section{The function}\label{sec:fn}

Let $\chi$ be the character modulo $5$ with
$\chi(0),\ldots,\chi(4) = 0,\,1,\,i,\,-i,\,-1$, primitive and odd. Put
\begin{equation}\label{eq:f}
f(s) \;=\; \tfrac12\bigl[(1-i\kappa)L(s,\chi) + (1+i\kappa)L(s,\bar\chi)\bigr].
\end{equation}
Davenport and Heilbronn~\cite{DH36} introduced this function and showed it has zeros off
the critical line; Titchmarsh~\cite[\S10.25]{Titchmarsh} is the standard exposition.

\paragraph{$\kappa$ from the Gauss sum \tier{T1}.} Write
$\varepsilon_\chi = \tau(\chi)/(i\sqrt5)$, so that
$\Lambda(s,\chi) = \varepsilon_\chi\Lambda(1-s,\bar\chi)$ with
$\Lambda(s,\chi) = (5/\pi)^{(s+1)/2}\Gamma(\tfrac{s+1}{2})L(s,\chi)$. Substituting
\eqref{eq:f} into $\Lambda_f(s) = (5/\pi)^{(s+1)/2}\Gamma(\tfrac{s+1}{2})f(s)$ and
demanding $\Lambda_f(s) = \Lambda_f(1-s)$, the coefficient of $\Lambda(s,\chi)$ matches
precisely when
\begin{equation}\label{eq:rootnumber}
(1-i\kappa)\,\varepsilon_\chi \;=\; 1+i\kappa,
\qquad\text{i.e.}\qquad \kappa = \tan\bigl(\tfrac12\arg\varepsilon_\chi\bigr),
\end{equation}
and the coefficient of $\Lambda(s,\bar\chi)$ then matches automatically because
$\varepsilon_{\bar\chi} = \overline{\varepsilon_\chi}$. Note that \eqref{eq:f} is
self-reciprocal while the equation for $L(s,\chi)$ alone exchanges $\chi$ and $\bar\chi$;
that exchange is exactly what \eqref{eq:f} is built to undo. Evaluating,
$\tan\arg\varepsilon_\chi = \sin 36^\circ/\sin 72^\circ = 1/\varphi$ with $\varphi$ the
golden ratio, whence the closed form
\[
\kappa \;=\; \frac{\sqrt{10-2\sqrt5}-2}{\sqrt5-1}
\;=\; 0.28407904384041229603\ldots
\]
Our $\kappa$ is computed as $\tan(\tfrac12\arg\varepsilon_\chi)$ from the Gauss sum, not
imported; it agrees with the closed form to $3.9\times10^{-62}$, and the residual of
\eqref{eq:rootnumber} is exactly $0$ at the working precision. The Dirichlet coefficients
are $(a_1,\ldots,a_5) = (1,\kappa,-\kappa,-1,0)$, real and periodic but \emph{not}
multiplicative: $a_2a_3 = -\kappa^2$ while $a_6 = a_1 = 1$.

\paragraph{Functional equation, measured \tier{T2}.} At twenty reproducibly chosen points
avoiding the critical line, the real axis and the gamma poles, the scale-aware relative
residual of $\Lambda_f(s) = \Lambda_f(1-s)$ is at worst $6.951\times10^{-58}$ at $60$
working digits and $8.092\times10^{-99}$ at $100$; it tracks the precision, which is the
check that it is a precision floor and not a defect. The relative form is necessary: one
of the twenty points has $|\Lambda_f| \approx 2\times10^{-31}$, where an absolute
residual would be meaningless.

\paragraph{Strip \tier{T1}.} For $\sigma \ge 2$,
$|f(s)-1| \le \sum_{n\ge2}|a_n|n^{-\sigma} \le \zeta(2)-1 = 0.644934066848\ldots < 1$.
Sharpening the same elementary estimate with the actual $|a_n|$ gives the abscissa
$\sigma_0$ with $\sum_{n\ge2}|a_n|n^{-\sigma_0} = 1$, namely
$\sigma_0 = 1.395136158235$, so every zero of $\Lambda_f$ lies in
$1-\sigma_0 < \operatorname{Re}s < \sigma_0$. The strip is wide --- this function has
zeros with $\operatorname{Re}s>1$ --- and it is what makes the census contour of \S3
meaningful without an arbitrary cutoff.

\paragraph{Coefficients of $-f'/f$ \tier{T1} and \tier{T2}.} In $\operatorname{Re}s>2$ set
$-f'(s)/f(s) = \sum_{n\ge2}b_nn^{-s}$, so that
$b_n = a_n\log n - \sum_{d\mid n,\,2\le d<n}b_d a_{n/d}$ with $b_1=0$. Exactly,
$b_2 = \kappa\log2$, $b_3 = -\kappa\log3$, and
\[
b_6 \;=\; (1+\kappa^2)\log 6 \;=\; 1.9363560766210387863,
\]
nonzero at a non-prime-power. For $\zeta$ the set of $n<50$ that are not prime powers and
have $b_n\ne0$ is empty; here it has $15$ members. The missing Euler product is therefore
visible in the coefficients the detector consumes, not only in the zeros. A second
computation of $b_n$ through the Dirichlet inverse of $(a_n)$ agrees with the recursion to
$1.245\times10^{-60}$ for $n\le200$.

\section{The window form}\label{sec:form}

Let $g$ be real and supported in $[-L,L]$, and set
\[
G(z) = \int g(u)e^{zu}\,du,\qquad h(v) = \int g(u+v)g(u)\,du,\qquad H(z) = G(z)G(-z),
\]
so $h$ is supported in $[-2L,2L]$ and $H(it) = |G(it)|^2$. The spectral form runs over
the zeros of the \emph{entire completed} function,
$Q_f(g) = \sum_\rho \operatorname{mult}(\rho)H(\rho-\tfrac12)$, and all reported values
are normalised as $Q_f(g)/\|g\|_2^2$. Moving the contour as in the companion paper's
\S4--5, with the conductor $5$ and the odd archimedean factor $\psi(\tfrac34+\tfrac{it}2)$
in place of the $\zeta$ data, gives
\begin{equation}\label{eq:form}
Q_f(g) \;=\; h(0)\log\frac{5}{\pi}
\;+\; \frac1{2\pi}\int_{\mathbb R}|G(it)|^2\operatorname{Re}\psi\!\left(\tfrac34+\tfrac{it}2\right)dt
\;-\; 2\!\!\sum_{2\le n\le e^{2L}}\!\! \frac{b_n\,h(\log n)}{\sqrt n}.
\end{equation}
There is \textbf{no pole term}: $\Lambda_f$ is entire, the trivial zeros of $f$ at
$s=-1,-3,\ldots$ cancelling the poles of the gamma factor. The arithmetic sum is
\textbf{finite} because $h$ vanishes beyond $2L$; at $L=2.6$ the cutoff is $n\le181$ with
$122$ nonzero terms, and the bound is decided in integer arithmetic and then re-tested as
$\log n < 2L$ at full precision, never by a binary float. The formula is of the
Weil--Yoshida--Bombieri type~\cite{Yoshida,Bombieri}; what is specific to $f$ is that
$b_n$ is not supported on prime powers.

\paragraph{Zero contributions.} An on-line conjugate pair $\{\tfrac12\pm i\gamma\}$
contributes $2H(i\gamma) = 2|G(i\gamma)|^2 \ge 0$. The complete orbit of a non-real
off-line zero, $O = \{\beta\pm i\gamma,\,1-\beta\pm i\gamma\}$, contributes, with
$\delta = \beta-\tfrac12$,
\begin{equation}\label{eq:orbit}
Q_O(g) \;=\; 4\operatorname{Re}\bigl(A\overline{B}\bigr) \;=\; 4\bigl(|P|^2-|S|^2\bigr),
\qquad A = G(\delta+i\gamma),\ B = G(-\delta+i\gamma),
\end{equation}
where $P = \int g\cosh(\delta u)e^{i\gamma u}du$ and
$S = \int g\sinh(\delta u)e^{i\gamma u}du$. This is a \emph{signed} rank-two form, with a
positive $\cosh$ channel and a negative $\sinh$ channel, and it is not
$|G(\delta+i\gamma)|^2$. As $\delta\to0$ it tends to $4H(i\gamma)$, the contribution of
\emph{two double} on-line zeros, so multiplicity is preserved: four zeros in, four out.

\paragraph{Gate A \tier{T2-N}.} The arithmetic side \eqref{eq:form} is compared with two
independent evaluations and three exact-algebra controls. The contour route writes
$Q_f(g) = (2/\pi)\operatorname{Re}\int_0^\infty(\Lambda_f'/\Lambda_f)(\sigma_c+it)
H(\sigma_c-\tfrac12+it)\,dt$ at $\sigma_c = 1.42$, to the right of every zero, with
$\Lambda_f'/\Lambda_f$ built from $f$ and $f'$ directly and \emph{never} from the $b_n$;
the zero route sums over the census of \S3.

\begin{center}
\begin{tabular}{@{}llll@{}}
\toprule
Check & Discrepancy & Declared budget & Tier \\
\midrule
arithmetic vs.\ logarithmic-derivative contour & $1.73\times10^{-13}$ & $1.316\times10^{-6}$ & T2-N \\
arithmetic vs.\ zero-side sum over the census & $5.65\times10^{-11}$ & $7.13\times10^{-5}$ & T2-N \\
control: negate the selected quartet (sign flip) & $1.83\times10^{-62}$ & exact algebra & T1 \\
control: omit $b_6$ at a window containing $n=6$ & $1.94\times10^{-61}$ & exact algebra & T1 \\
control: flip each arithmetic term in turn & $3.11\times10^{-61}$ & exact algebra & T1 \\
\bottomrule
\end{tabular}
\end{center}

\noindent The quartet control negates the selected orbit's contribution on the zero side
alone and checks that the discrepancy moves by exactly twice it; with
$Q_{O_1}(g) = 8.720121523693614\times10^{-11}$ on that low-frequency test function, the
predicted and observed changes agree to $1.83\times10^{-62}$. The $b_6$ control predicts
a change of $+2b_6h(\log6)/\sqrt6 = -0.379255712825982$ from
$b_6 = 1.93635607662104$ and $h(\log 6) = -0.239879170384883$, and observes it to
$1.94\times10^{-61}$. The term-level sign flip negates each of the archimedean,
conductor and arithmetic terms in turn; the conductor row is the control that would catch
a wrong conductor, and it reproduces $-2\log5 = -3.21887582486820$ for a unit-normalised
test function in every sector and window. \textbf{Classification: PASS-N} --- converged
numerical agreement, explicitly uncertified, because the quadrature in the contour route
and the zero density in the zero-side tail are a library estimate and a measured quantity,
not rigorous bounds. The three controls, which are what would catch a wrong factor, a
wrong conductor or a dropped term, are exact.

\section{The zero census to $T=200$}\label{sec:census}

Counting by the argument principle on the strip $-0.4<\operatorname{Re}s<1.4$, whose
vertical edges are provably zero-free by \S\ref{sec:fn}, gives $130$ zeros with
$0<\operatorname{Im}s<200$ against a Riemann--von Mangoldt main term of $129.548$; the
worst deviation of a winding number from an integer over the whole ladder is
$1.32\times10^{-23}$. Sign changes of the Hardy-type $Z$ give $122$ critical-line
ordinates, and four off-line zeros with $\operatorname{Re}s>\tfrac12$ are located by box
subdivision and Newton, each carrying local winding number exactly $1$ on a circle of
radius $0.02$, hence simple. The reconciliation
\[
122 \;+\; 2\cdot4 \;=\; 130
\]
is exact, and it is also what closes the band $\tfrac12<\operatorname{Re}s<0.55$ that the
located search deliberately avoids: an unlocated zero there would break it.

%% generated by note/make_tables.py from the data files
%% DO NOT EDIT BY HAND
\begin{center}\small
\begin{tabular}{@{}llllll@{}}
\toprule
 & $\beta$ & $\gamma$ & $\delta=\beta-\tfrac12$ & $|f(\rho)|$ & winding \\
\midrule
$O_1$ & $0.808517182456637386$ & $85.6993484853775922$ & $0.308517182456637386$ & $1.75\times10^{-60}$ & $1$ \\
$O_2$ & $0.650830080609737082$ & $114.163342730756981$ & $0.150830080609737082$ & $3.21\times10^{-60}$ & $1$ \\
$O_3$ & $0.574356050450805991$ & $166.479305913168156$ & $0.0743560504508059907$ & $5.99\times10^{-60}$ & $1$ \\
$O_4$ & $0.724257694626809780$ & $176.702461242855825$ & $0.224257694626809780$ & $2.99\times10^{-60}$ & $1$ \\
\bottomrule
\end{tabular}
\end{center}


\noindent The residuals are at the working precision of the census run ($60$ digits); at
$140$ digits the refined $O_1$ has $|f(\rho_1)| = 8.2\times10^{-141}$ with Newton step
$6.5\times10^{-141}$ and is stable to $3.17\times10^{-106}$ between $100$ and $140$
digits. These are roots to the precision floor, not approximate ones.

\paragraph{Which are cross-checked \tier{T2}.} \textbf{One of the four, and only against
a secondary source.} $O_1$ was found from a decimal seed supplied with the task; Newton
moved that seed by $7.8\times10^{-15}$, so the seed's digits are correct to about
fourteen places. The provenance of that decimal is a secondary report of the literature
computations of Spira~\cite{Spira} and Balanzario and
S\'anchez-Ortiz~\cite{BSO}; \textbf{neither paper was read} --- the AMS PDF of the latter
returned HTTP~403 from the machine used --- and no numerical value from either appears
in this note. $O_2$, $O_3$ and $O_4$ are \textbf{not} cross-checked against anything: they
are this computation's own. Bombieri and Ghosh~\cite{BG11} is the standard modern
reference for the analytic theory of $f$; it was not read either, and no result or
constant of theirs is used --- in particular the counting function above is the computed
argument-principle count, with the main term quoted only as a diagnostic.

\section{The certified witnesses}\label{sec:witness}

\tier{T2-C} Two test functions are frozen, with every coefficient stored as an exact
$30$-digit decimal string. Both are odd, both are in the span of a low block together
with a block resonant at $\gamma_1$, and the Gram matrix of the normalised basis is the
identity exactly (deviation $4.83\times10^{-60}$ and $7.4\times10^{-60}$), so the
generalised eigenproblem is the ordinary one and no direction was regularised or
discarded. The coefficients of $g_*$ are in Appendix~A.

\begin{center}
\begin{tabular}{@{}lll@{}}
\toprule
& $g_*$ & $g_\sharp$ \\
\midrule
$L$ & $2.0$ & $1.8$ \\
sector & odd & odd \\
$N$ & $37$ & $57$ \\
arithmetic cutoff & $n\le56$, $39$ nonzero terms & $n\le38$, $26$ nonzero terms \\[2pt]
$Q_f/\|g\|_2^2$ & $-0.679071346441480590986$ & $-1.68123895523670103\times10^{-16}$ \\
ball radius & $1.39\times10^{-116}$ & $8.15\times10^{-177}$ \\
strictly negative & yes & yes \\
\bottomrule
\end{tabular}
\end{center}

\paragraph{Why the certificate is complete.} The enclosure is of the \emph{whole}
quantity, not of an approximation to it, and there is nothing outside the ball:

\begin{itemize}\itemsep2pt
\item \textbf{The arithmetic sum is finite and complete.} $h$ vanishes beyond $2L$, so
\eqref{eq:form} has $39$ terms at $L=2$ and no tail is dropped.
\item \textbf{No quadrature.} The archimedean term is the closed form in digamma, the
Hurwitz zeta and the Lerch transcendent; Arb encloses all three.
\item \textbf{No truncation and no spectral input.} $Q_f(g_*)$ is computed from the
arithmetic side alone. It does not know where any zero is, so the census plays no part in
it and the sign does not depend on the census being right.
\item \textbf{The test function carries no error.} Its coefficients are exact finite
decimals; the ball encloses $Q_f$ for \emph{this} $g_*$, not for a rounding of some other
one.
\item \textbf{Independent implementation.} The certifying code re-implements the
cross-correlations, $\kappa$, the $b_n$ and the archimedean closed form against Arb,
sharing no code with the \texttt{mpmath} assembly; the two agree to all $44$ digits
\texttt{mpmath} carries, so the certificate is simultaneously a cross-check of both.
\end{itemize}

\noindent The only inputs are Arb's enclosures and the derivation of \eqref{eq:form}.
$\kappa$ enters through square roots only. This is \tier{T2-C} and not a formal
verification.

\section{The decomposition}\label{sec:decomp}

\tier{T2-C} Both terms below are finite sums and both are enclosed. At $g_*$:
\begin{equation}\label{eq:decomp}
\underbrace{-0.679071346441480590986}_{Q_f(g_*)}
\;=\;
\underbrace{-0.985198326439305481794}_{Q_{O_1}(g_*)}
\;+\;
\underbrace{+0.306126979997824890807}_{\Qrest(g_*)},
\end{equation}
with ball radii $1.39\times10^{-116}$, $4.06\times10^{-117}$ and $1.63\times10^{-116}$,
and additive reconciliation residual $2.99\times10^{-116}$. Here $\Qrest = Q_f - Q_{O_1}$
is evaluated by subtracting the exact rank-two orbit form \eqref{eq:orbit} from the
arithmetic side, so \emph{every} remaining zero --- including every off-line zero above
the census height and every zero never located --- is included implicitly. Nothing is
replaced by a finite zero list.

The selected quartet therefore carries more than the whole of the negativity, and the
rest of the spectrum returns $31.07\%$ of it. At $g_\sharp$ the same decomposition reads
$-6.45917960933013\times10^{-14} + 6.44236721977777\times10^{-14}$: near the onset the
cancellation is $99.7397\%$, which is why the witness there is of order $10^{-16}$ rather than
$10^{-14}$.

\paragraph{The remaining-zero form is positive on every space searched \tier{T2-N}.}
$\lmin$ of the restriction of $\Qrest$ to the level-$2$ space, alongside $\lmin(Q_f)$:

%% generated by note/make_tables.py from the data files
%% DO NOT EDIT BY HAND
\begin{center}\small
\begin{tabular}{@{}lllll@{}}
\toprule
$L$ & $\lmin(Q_f)$, even & $\lmin(\Qrest)$, even & $\lmin(Q_f)$, odd & $\lmin(\Qrest)$, odd \\
\midrule
$1.0$ & $4.034840454024\times10^{-5}$ & $4.03625203581\times10^{-5}$ & $0.01356156739650$ & $0.0132281100537$ \\
$1.2$ & $8.780065301754\times10^{-9}$ & $8.46979844459\times10^{-9}$ & $9.702616267333\times10^{-6}$ & $9.54583799747\times10^{-6}$ \\
$1.4$ & $5.087177960475\times10^{-14}$ & $4.64125898422\times10^{-14}$ & $6.538829304189\times10^{-11}$ & $6.63381728826\times10^{-11}$ \\
$1.6$ & $2.549827775205\times10^{-19}$ & $2.25086996550\times10^{-19}$ & $2.153851664770\times10^{-16}$ & $2.21332806419\times10^{-16}$ \\
$1.8$ & $7.580609464352\times10^{-24}$ & $7.01854463036\times10^{-24}$ & $1.007541951495\times10^{-20}$ & $1.08891994498\times10^{-20}$ \\
$2.0$ & $-0.4896520626161$ & $2.02946591361\times10^{-27}$ & $-0.6790713464415$ & $1.20623694960\times10^{-24}$ \\
$2.2$ & $-1.219875374653$ & $1.22467880540\times10^{-30}$ & $-1.178435305399$ & $4.87721021342\times10^{-28}$ \\
$2.4$ & $-1.503014486764$ & $1.81134931584\times10^{-33}$ & $-1.590268954821$ & $4.26137491134\times10^{-31}$ \\
$2.6$ & $-2.137655567743$ & $8.47226605297\times10^{-36}$ & $-2.005507257484$ & $8.13613127375\times10^{-34}$ \\
\bottomrule
\end{tabular}
\end{center}


\noindent All $18$ rows are positive. Once $Q_f$ is macroscopically negative the
remaining-zero form does not merely stay positive: it collapses towards zero, reaching
$8.47\times10^{-36}$ at $L=2.6$ while $Q_f$ is at $-2.14$. Removing $O_1$ from the frozen
witness's space and re-optimising gives $\lmin = +1.20623694960\times10^{-24}$; no other
negative direction is found there, which is a statement about that space and not about
positivity.

\paragraph{The consequence, stated exactly.} The remaining-zero theorem of the companion
paper assumes a remaining-zero form that is nonnegative, and states that assumption as
\emph{sufficient}, not necessary, for its conclusion. On this object that hypothesis
\textbf{holds on every space searched}, so the theorem's conclusion applies here; and
because its proof uses the hypothesis only through the single inequality
$\Qrest \ge 0$, nothing else about where the other zeros lie is needed. The masking
contribution is $31.07\%$ in the witness direction and does \emph{not} reverse the sign.
What is \emph{not} established is that a positive finite restriction establishes the
whole-window hypothesis: controlling the infinite-dimensional complement and its coupling
to the block would be needed, and a positive Galerkin matrix supplies neither.

\paragraph{And the confound, once more.} $\Qrest$ is built from the on-line zeros of an
object with no Euler product, and it is positive; that is the strongest evidence
\emph{internal to this object} that the negativity is attributable to the off-line
quartets, and it is not a proof.

\section{Onset}\label{sec:onset}

\tier{T2-N} The scan runs $L = 1.0,1.2,\ldots,2.6$ in both parity sectors at three
declared basis levels; enrichment can only lower $\lmin$, and the levels are reported as
enrichments, never folded into one another. Over all $54$ rows the Rayleigh quotient of
the returned eigenvector agrees with the eigenvalue to $5.76\times10^{-60}$, so the
eigensolver is confirmed by direct evaluation. At the deepest level, $\lmin$ falls by a
factor $1.459\times10^{17}$ from $L=1.0$ to $L=1.6$ while remaining positive, and then
changes sign:

%% generated by note/make_tables.py from the data files
%% DO NOT EDIT BY HAND
\begin{center}\footnotesize
\begin{tabular}{@{}lllllll@{}}
\toprule
$L$ & even, lvl 1 & even, lvl 2 & even, lvl 3 & odd, lvl 1 & odd, lvl 2 & odd, lvl 3 \\
\midrule
$1.6$ & $1.97261814306\times10^{-13}$ & $2.54982777520\times10^{-19}$ & $2.66141166363\times10^{-22}$ & $2.91871108168\times10^{-11}$ & $2.15385166477\times10^{-16}$ & $1.09593824675\times10^{-18}$ \\
$1.8$ & $2.91410423267\times10^{-15}$ & $7.58060946435\times10^{-24}$ & $-2.66180729007\times10^{-21}$ & $1.19062366907\times10^{-13}$ & $1.00754195149\times10^{-20}$ & $-1.68123895524\times10^{-16}$ \\
$2.0$ & $-0.188728484539$ & $-0.489652062616$ & $-0.547390086184$ & $-0.658652328928$ & $-0.679071346441$ & $-0.710175512230$ \\
\bottomrule
\end{tabular}
\end{center}


\noindent Twenty-six of the $54$ rows are negative. The crossing is not gradual: the form
sits on a knife edge through most of the scanned range and then falls off it, from
$-1.68\times10^{-16}$ at $L=1.8$ to $-0.71$ at $L=2.0$.

\paragraph{Which sector detects first.} \textbf{The scan does not resolve this.} Both
sectors are positive at $L=1.6$ and negative at $L=1.8$ at level $3$, so on this grid they
cross together; the odd value is deeper by a factor $6.316\times10^{4}$ at $L=1.8$ and
by $0.16$ at $L=2.0$. Since $K_{\mathrm{even}}(2.0,\gamma_1)/K_{\mathrm{odd}}(2.0,\gamma_1)
= 1.00634$ at this height, the two sectors have essentially the same isolated
sensitivity, and \textbf{no structural explanation is offered} for the depth difference.

\paragraph{Two comparisons.} Both are computed by the verifier from the Phase-A values of
$(\delta_1,\gamma_1)$, not quoted from elsewhere.

\emph{(i) Bandwidth scale.} With
$L_{\mathrm{pred}} = \tfrac12\log(5\gamma_1/2\pi) = 2.11120303$, and the measured onset
bracketed below by the certified negative witness at $L=1.8$ and above by the macroscopic
one at $L=2.0$,
\[
\frac{L_{\mathrm{pred}}}{1.8} = 1.17289, \qquad
\frac{L_{\mathrm{pred}}}{2.0} = 1.0556 .
\]
This is a reference scale, not a predicted onset: its provenance is a bandwidth heuristic
that was not verified, and the bracket is one-sided in the sense of non-claim~5.

\emph{(ii) The $\delta^2$ coefficient.} In the odd sector the exact extremal constant is
$K_{\mathrm{odd}}(L,\gamma) = \int_{-L}^{L}u^2\cos^2(\gamma u)\,du = 2.6582355735990895$
at $L=2$, $\gamma=\gamma_1$, and the companion paper's $\delta^2$ coefficient gives
\[
4K_{\mathrm{odd}}(2.0,\gamma_1)\,\delta_1^2 \;=\; 1.01207377,
\qquad\text{against}\qquad |Q_{O_1}(g_*)| = 0.985198326439305,
\]
a ratio of $1.02728$. Measured instead against the actual drop from the collapsed
configuration, $Q_{O_1}(0) - Q_{O_1}(\delta_1) = 0.986028353$ with
$Q_{O_1}(0) = 4|f\,\hat{}\,(\gamma_1)|^2 = 0.000830026353$, the ratio is $1.02641$.

\emph{Is that inside what the remainder allows?} Yes, and \emph{so comfortably that the
comparison is not a sharp test of the bound}. At $L\delta_1 = 0.617034$ the companion
paper's $\delta^4$ allowance is
$\bigl(\tfrac25+\tfrac{16}3e^{2L\delta}\bigr)L^5\delta^4 = 5.42744$, while the observed
discrepancy $|4K\delta^2 - \text{drop}|$ is $0.02605$ --- smaller by a factor $208.4$. The
$2.6\%$ agreement is therefore \emph{better} than the bound requires, by more than two
orders of magnitude, and should not be read as confirming the bound's constant; the bound
is simply very loose at this displacement. $L\delta_1 = 0.617$ is far outside any
small-displacement regime, and no small-$\delta$ law is used anywhere in \S4 or \S5.

\section{What this rules out, and what it does not}\label{sec:rules}

\paragraph{Ruled out.} Any claim that finite-window Weil positivity follows from a
functional equation of $\zeta$'s shape, together with entireness of the completion and a
zero-containing strip. The Davenport--Heilbronn function has all three --- \eqref{eq:f}
is self-reciprocal, $\Lambda_f$ is entire, and every zero lies in
$|\operatorname{Re}s-\tfrac12| < 0.895136$ --- and its window form at $L=2.0$ is negative, at $-0.679071346441$, now by a \emph{certified example at a stated window} rather than by citation of
the existence of off-line zeros. The same example rules out the weaker claim that a
positive finite Galerkin restriction of a remaining-zero form establishes a whole-window
hypothesis: here $\lmin(\Qrest)>0$ on every space searched while $Q_f$ is macroscopically
negative at five of the nine windows.

\paragraph{Not ruled out, and not claimed.} Nothing about $\zeta$: no statement here has a
$\zeta$ analogue, and the whole point of the control is that it runs where positivity is
known to fail. Nothing that separates the two properties of $f$, off-line zeros and the
absent Euler product, which is the confound of non-claim~3. Nothing about other
$L$-functions, other conductors, other bases, or other windows than the nine scanned.
Nothing about positivity at $L\le1.6$: those are non-detections. And nothing certified
about the \emph{onset}, only about the two witnesses --- certifying the crossing would
need a certified lower bound at $L=1.6$ over a whole space, which is a different and
harder object than a certified value at a point.

\paragraph{Open, in four sentences.} Whether $\lmin(\Qrest)$, already at
$8.47\times10^{-36}$ at $L=2.6$, stays positive as the basis grows, or whether its
near-null direction is the leading edge of a genuine negative one. Whether the onset can
be bracketed by certified bounds at both ends rather than by a certified negative and a
non-detection. Whether the census can be made rigorously complete, the winding numbers
being integers to $1.32\times10^{-23}$ by adaptive sampling rather than by interval
arithmetic. And whether an off-line quartet of small $\delta$ at larger height --- $O_3$
has $\delta = 0.0743560504508059907$ at $\gamma = 166.479305913168156$ --- is seen earlier
or later than $O_1$ by a basis resonant at \emph{its} height, which the present scan,
resonant only at $\gamma_1$, cannot say.

\appendix

\section{The frozen witness $g_*$}\label{app:witness}

$L = 2.0$, odd sector, $N = 37$. The test function is
$g_*(u) = \sum_{m} c_m\,\phi_{k_m}(u)$ with
$\phi_k(u) = \sin\bigl((k+1)\pi u/L\bigr)\big/\sqrt{L}$ on $[-L,L]$ and zero outside, so
that $\|g_*\|_2^2 = \sum_m c_m^2$. The index set is the low block $k = 0,\ldots,15$
together with the resonant block $k = 44,\ldots,64$, chosen so that
$w_k = (k+1)\pi/L$ brackets $\gamma_1$. The coefficients are exact finite decimals, read
back from the certificate as strings; the table below is generated mechanically from
\src{certificates/witness.json} and is regenerated and compared by the verifier.

%% generated by note/make_witness_table.py from certificates/witness.json
%% DO NOT EDIT BY HAND
\begin{center}\footnotesize
\begin{tabular}{@{}rrlrrl@{}}
\toprule
$m$ & $k_m$ & $c_m$ & $m$ & $k_m$ & $c_m$ \\
\midrule
$0$ & $0$ & $-0.0026088367880077501762848431774$ & $19$ & $47$ & $0.0160794083499821171953260941976$ \\
$1$ & $1$ & $0.0059367297983683678676143712306$ & $20$ & $48$ & $0.00498986780562686767195561151428$ \\
$2$ & $2$ & $-0.0185534062780231023921261896105$ & $21$ & $49$ & $0.0000277053929521654177814882586319$ \\
$3$ & $3$ & $0.00544846305151052127579958739366$ & $22$ & $50$ & $-0.0222179267261014551646436764475$ \\
$4$ & $4$ & $-0.0156396017373944635966957109136$ & $23$ & $51$ & $0.039586845834277726764109564134$ \\
$5$ & $5$ & $-0.0109783164655112390655789188025$ & $24$ & $52$ & $0.0942302999106238255858728047973$ \\
$6$ & $6$ & $-0.0125907754313286158626601747424$ & $25$ & $53$ & $0.804187411949596545617290494929$ \\
$7$ & $7$ & $-0.0144628048949237093476726836875$ & $26$ & $54$ & $-0.559992992803302061132173856712$ \\
$8$ & $8$ & $-0.0146343914894490798881726228996$ & $27$ & $55$ & $0.137708858583258544659438566561$ \\
$9$ & $9$ & $0.0103898524391239437098974349874$ & $28$ & $56$ & $0.0221776726674697424766634785707$ \\
$10$ & $10$ & $0.00448607789902456302212489057999$ & $29$ & $57$ & $0.0150454955403207338083794667482$ \\
$11$ & $11$ & $0.0241973863830400609309747239095$ & $30$ & $58$ & $-0.0303098024499984095981337680134$ \\
$12$ & $12$ & $-0.0141545580712512199946130198503$ & $31$ & $59$ & $0.0273927617083722981863389567813$ \\
$13$ & $13$ & $0.00838291335579069185266302436363$ & $32$ & $60$ & $0.00112786440916397755739920517311$ \\
$14$ & $14$ & $0.00306472453668078854500166029552$ & $33$ & $61$ & $-0.0226289801121897281807111788614$ \\
$15$ & $15$ & $0.0319718989223066404550495907318$ & $34$ & $62$ & $0.00211860269767553369614679144869$ \\
$16$ & $44$ & $-0.0311893304913935554292879513548$ & $35$ & $63$ & $0.019589441263546411704813632605$ \\
$17$ & $45$ & $-0.00554065734427570197610267400465$ & $36$ & $64$ & $0.030112760581951388414172362309$ \\
$18$ & $46$ & $-0.0304191101621335069943242266561$ & & & \\
\bottomrule
\end{tabular}
\end{center}


\noindent The mass is concentrated on $k=53$ and $k=54$, whose frequencies bracket
$\gamma_1 = 85.699348485377592$; the witness has normalised overlap $0.9513529401$ with
$u\cos(\gamma_1 u)$, the extremiser of the isolated derivative-evaluation functional,
while the ground state of $\Qrest$ is orthogonal to it to $9.07\times10^{-12}$.

\section{Retraction log}\label{app:retr}

Six corrections were made during this work. They are reproduced from
\src{RETRACTIONS.md}, in order.

\begin{description}\itemsep3pt
\item[R1. Two-sided zero-count normalisation.] The Riemann--von Mangoldt main term was
written $(T/\pi)\log(5T/2\pi e)$, which counts $|\operatorname{Im}\rho|<T$, while the
census counts $0<\operatorname{Im}\rho<T$. Corrected to $(T/2\pi)\log(5T/2\pi e)$. The
near-exact factor of two between the winding numbers and the ``main term'' is what
exposed it. Diagnostic only: the counts come from the argument principle, and no reported
number changed.

\item[R2. Off-line search contour too close to the critical line.] The first
implementation subdivided from $\operatorname{Re}s \ge 0.505$, which passes within
$0.005$ of the on-line zeros; $\arg f$ then swings by nearly $\pi$ between samples and the
adaptive refinement recursed to its depth limit on every band. Corrected by searching at
$\operatorname{Re}s\ge0.55$ and closing the uncovered band by the reconciliation of \S3,
which is a stronger check and cheap. Method only; no reported number changed.

\item[R3. Division by zero at a removable singularity.] The closed form for
$\int\cos(wu)e^{zu}du$ has a $0/0$ at $z = \pm iw$. A unit test evaluating there caught
it; an explicit limit branch was added. The condition never arose in a reported
computation --- $H$ is evaluated at zero ordinates, not at basis frequencies --- and the
branch exists so the test can cover the case.

\item[R4. The first independent archimedean route was too weak to be a check.] It
evaluated the frequency-side integral by quadrature; the integrand oscillates with period
$\pi/L$ and the result was accurate only to about $10^{-6}$, which cannot distinguish a
correct assembly from one with a small error. Replaced by Gauss's integral for $\psi$,
giving a one-dimensional non-oscillatory integral with an exact tail. The two routes now
agree to $6.225\times10^{-21}$ against an analytic bound of $1.66\times10^{-20}$ on the
piece deliberately omitted near the origin.

\item[R5. A gate threshold that could not be met.] Phase B required agreement below
$10^{-40}$ while its own declared budget was $10^{-20}$. A gate that cannot pass is not a
gate; the threshold was tied to the route's own budget. No reported value changed.

\item[R6. Decimal window lengths parsed at the ambient precision.] \textbf{This is the one
that mattered.} \texttt{mpmath}'s \texttt{mpf} rounds a decimal string to the
\emph{current} working precision and freezes it. A window length built as
\texttt{mpf('1.8')} before the precision was raised carried only $15$ digits, and the
subsequent precision ladder at $60$, $100$ and $160$ digits then evaluated the form at
that same slightly-wrong $L$ three times, agreed with itself to $10^{-85}$, and reported
a value wrong in the \emph{thirteenth} significant digit. A ladder that agrees with
itself to $85$ digits is not evidence that the number is right. It was caught only
because two scripts froze the same witness independently and disagreed in the $13$th
digit --- the ladder itself could not have caught it, and would not have, at any
precision. Every decimal constant that feeds the form now goes through a parser fixed at
$220$ digits, and a guard test fails if any of them is reconstructed at the ambient
precision. Blast radius: the certified witness sits at $L=2.0$, exactly representable in
binary, and is unaffected; the window scan ran at a consistent precision throughout and is
unaffected; the $L=1.8$ witness was recomputed, and its sign and first twelve digits are
unchanged.
\end{description}

\section{The two literature witnesses}\label{app:lit}

Two contemporaneous documents report negative Davenport--Heilbronn window witnesses.
\textbf{Only their landing pages were read}; neither PDF was retrieved, and nothing below
is a reading of their arguments.

The first is T.~Ziad, \emph{Localized Constrained Weil Forms: Near-Degeneracy and a
Davenport--Heilbronn First Kernel}, Zenodo,
\href{https://doi.org/10.5281/zenodo.23145455}{doi:10.5281/zenodo.23145455}, 7 October
2026, which reports ``an exactly constrained negative witness at $a = 9/5$'' for odd,
\emph{constrained} test functions, and a first loss-of-positivity scale strictly between
$\tfrac12\log2$ and $9/5$. The second is a Zenodo series on cross-energy of Dirichlet
partial sums, from which the sentence ``positive up to a window $L \approx 3.6$ \ldots
$\lambda_{\min} = -0.73$ at $L=4$, with a minimizer concentrated on its off-line zero
$0.8085+85.70i$'' was surfaced.

The two window numbers differ by exactly a factor of two. In the convention used here,
$g$ is supported in $[-L,L]$ and its autocorrelation $h$ in $[-2L,2L]$, so a factor of two
is what one would expect between a document measuring the support of $g$ and one
measuring the support of $h$. \textbf{That is a hypothesis, not a finding}: nothing here
establishes which convention either document uses, and until it is settled neither
agreement nor disagreement with them can be asserted. Our level-$3$ value at $L=2.0$,
$-0.7101755122300$, is in the same region as the quoted $-0.73$ under the hypothesis, and
is an upper bound for the true minimum because it is computed on a finite subspace --- but
that observation is worth exactly as much as the convention hypothesis it rests on.

\textbf{Nothing in this note was fitted to either value.} The scanned window range was
fixed before either document was located, no threshold or basis was tuned, and no
novelty is claimed: at least two contemporaneous efforts address the same control, one of
them with its own interval certificate.

\section{Reproduction}\label{app:repro}

Code and data: \url{https://github.com/innerlightr-wq/weil-window-control}, directory
\src{extensions/davenport-heilbronn-2026-10-23/}. Python~3.14.7 (CPython),
\texttt{mpmath}~1.3.0, \texttt{python-flint}~0.9.0 (Arb). No \texttt{numpy}, no
\texttt{scipy}, and no \texttt{float64} on any path reaching a reported value. All
conventions, precisions, basis levels, windows and thresholds are in \src{src/config.py};
no randomness is used.

\begin{center}\small
\begin{tabular}{@{}llll@{}}
\toprule
Result & Script & Data & Runtime \\
\midrule
\S\ref{sec:fn} function, $\kappa$, functional equation & \src{src/phaseA.py} & \src{results/phaseA.json} & ${\approx}3$\,m \\
\S\ref{sec:fn} $b_n$, \S\ref{sec:form} form & \src{src/phaseB.py} & \src{data/dh\_coefficients.csv} & ${\approx}6$\,m \\
\S\ref{sec:form} orbit conventions & \src{tests/test\_orbits.py} & --- & ${\approx}2$\,m \\
\S\ref{sec:census} census & \src{src/census.py} & \src{data/dh\_zeros.csv} & ${\approx}25$\,m \\
\S\ref{sec:form} Gate A & \src{src/phaseC.py}, \src{src/gate\_signflip.py} & \src{data/gate\_checks.csv}, \src{data/gate\_signflip.csv} & ${\approx}25$\,m \\
\S\ref{sec:onset} scan & \src{src/phaseD.py} & \src{data/window\_scan.csv} & ${\approx}45$\,m \\
\S\ref{sec:witness} freeze & \src{src/freeze.py} & \src{certificates/witness*.json} & ${\approx}20$\,m \\
\S\ref{sec:witness} certificates & \src{src/certify.py} & \src{certificates/ball\_certificate*.txt} & ${\approx}2$\,m \\
\S\ref{sec:decomp} remaining form & \src{src/phaseE.py} & \src{data/remaining\_form.csv} & ${\approx}15$\,m \\
\S\ref{sec:decomp} masking & \src{src/phaseF.py} & \src{data/masking.csv} & ${\approx}5$\,m \\
\S\ref{sec:onset} $K$, perturbation & \src{src/phaseG.py} & \src{data/sensitivity.csv} & ${\approx}10$\,m \\
\bottomrule
\end{tabular}
\end{center}

\noindent The two certificates are reproduced by
\begin{verbatim}
.venv/bin/python src/certify.py 400
.venv/bin/python src/certify.py 600 certificates/witness_smallest_L.json \
                                    ball_certificate_L1.8.txt
\end{verbatim}

\noindent \texttt{make verify}, in the note directory, runs
\src{note/verify\_note\_numbers.py}: it rebuilds every quantity printed above from the
data files, scans this document for numeric literals, fails if any literal is not derived
or is printed to more significant digits than its source carries, regenerates
Appendix~A's table from \src{certificates/witness.json} and compares it, and cross-checks
the Gate~A values against \src{REPORT.md}. The full suite is \src{./run\_suite.sh}.

\section*{Licence}
\addcontentsline{toc}{section}{Licence}

As the parent repository: code MIT, this note and the data CC~BY~4.0. CC~BY~4.0 requires
attribution; if you use the note, the certificates or the data, please cite this note.

\section*{Acknowledgments}
\addcontentsline{toc}{section}{Acknowledgments}

AI assistance was used for algebra checking, numerical verification, adversarial review,
literature organization, and manuscript drafting. All conceptual development, research
direction, interpretation, and final responsibility for the content belong to the author.

\begin{thebibliography}{99}

\bibitem{BSO} E.~P. Balanzario and J.~S\'anchez-Ortiz,
\emph{Zeros of the Davenport--Heilbronn counterexample},
Math. Comp. \textbf{76} (2007), no.~260, 2045--2049.

\bibitem{Bombieri} E. Bombieri,
\emph{Remarks on Weil's quadratic functional in the theory of prime numbers, I},
Atti Accad. Naz. Lincei Cl. Sci. Fis. Mat. Natur. Rend. Lincei (9) Mat. Appl. \textbf{11}
(2000), no.~3, 183--233.

\bibitem{BG11} E. Bombieri and A. Ghosh,
\emph{Around the Davenport--Heilbronn function},
Russian Math. Surveys \textbf{66} (2011), no.~2, 221--270;
Uspekhi Mat. Nauk \textbf{66} (2011), no.~2, 15--66.

\bibitem{DH36} H. Davenport and H. Heilbronn,
\emph{On the zeros of certain Dirichlet series},
J. London Math. Soc. \textbf{11} (1936), 181--185; \textbf{11} (1936), 307--312.

\bibitem{DeJesus} E. De Jes\'us,
\emph{A $\delta^2$--$L^3$ Sensitivity Law for Finite-Window Weil Positivity, with a
Function-Field Control of the Connes--van Suijlekom Pipeline}, Zenodo, 2026.
\href{https://doi.org/10.5281/zenodo.23250932}{doi:10.5281/zenodo.23250932}.

\bibitem{Arb} F. Johansson,
\emph{Arb: Efficient arbitrary-precision midpoint-radius interval arithmetic},
IEEE Trans. Comput. \textbf{66} (2017), no.~8, 1281--1292.

\bibitem{Spira} R. Spira,
\emph{Some zeros of the Titchmarsh counterexample},
Math. Comp. \textbf{63} (1994), no.~208, 747--748.

\bibitem{Titchmarsh} E.~C. Titchmarsh,
\emph{The Theory of the Riemann Zeta-Function}, 2nd ed., revised by D.~R. Heath-Brown,
Oxford University Press, 1986, \S10.25.

\bibitem{Yoshida} H. Yoshida,
\emph{On Hermitian forms attached to zeta functions},
in: N. Kurokawa and T. Sunada (eds.), \emph{Zeta Functions in Geometry},
Adv. Stud. Pure Math. \textbf{21}, Math. Soc. Japan, Kinokuniya, Tokyo, 1992, 281--325.

\end{thebibliography}

\end{document}
